\section{The original target}

\textcite{bresnan2007} is the anchor because it treats the alternation as a
grammatical problem, not a simple frequency count: the claim isn't that
whatever occurs often is grammatical, but that observed choices expose
systematic constraints on what speakers treat as available. That matters because
informal judgements compress several diagnostics into one feeling. A
construction may sound odd because no viable analysis is available, because a
verb--construction pairing conflicts with an established value, because the
relevant discourse or communicative situation is missing, because speakers have
too little evidence to settle a rare form's status, because an established
competitor keeps winning where the form might otherwise appear, or because
processing pressure makes a recoverable structure feel bad. Corpus data,
especially when normalized to the relevant opportunity set and paired with
judgement evidence, can help sort these cases.

The empirical target is production choice. In Bresnan et al.'s data the response
is whether the recipient appears as a noun phrase (NP), in the double-object
pattern, or as a prepositional phrase (PP), in the prepositional dative, with
predictors for the verb, recipient, theme, discourse status, and corpus. This
paper preserves that target and treats the production probability as evidence to
interpret, not as a direct measure of grammaticality.
